

\documentclass[12pt]{beamer}
\usepackage{xcolor}
\usepackage{pgf,pgfarrows,pgfnodes,pgfautomata,pgfheaps,pgfshade}
\usetheme{Air}


\DeclareMathOperator*{\argmax}{arg\,max}

\title[CSC349A Numerical Analysis]{CSC349A Numerical Analysis}


\logo{\pgfputat{\pgfxy(-0.5,7.5)}{\pgfbox[center,base]{\includegraphics[width=1.0cm]{figures/uvic}}}}  
\beamertemplatenavigationsymbolsempty

    \defbeamertemplate{footline}{author and page number}{%
      \usebeamercolor[fg]{page number in head/foot}%
      \usebeamerfont{page number in head/foot}%
      \hspace{1em}\insertshortauthor\hfill%
      \insertpagenumber\,/\,\insertpresentationendpage\kern1em\vskip2pt%
    }
    \setbeamertemplate{footline}[author and page number]{}



\subtitle[Leture 4]{Lecture 4}
\date[2025]{2025}
\author[George Tzanetakis]{George Tzanetakis}
\institute[University of Victoria]{University of Victoria}
%\logo{\includegraphics[scale=.25]{unilogo.pdf}}


\usepackage{amsmath}
\usepackage{ulem}
\usepackage{color,soul}

\begin{document}
\frame{\maketitle} % <-- generate frame with title


\AtBeginSection[]
{
\begin{frame}<beamer>[allowframebreaks]{Table of Contents}
\tableofcontents[currentsection,currentsubsection, 
    hideothersubsections, 
    sectionstyle=show/shaded,
]
\end{frame}
}

\section{Number systems}

\begin{frame}{Round-off errors}

Round-off errors originate from two factors:

\begin{itemize}
\item{finite representations of possibly infinitely long numbers}
\item{finite range of values from a possibly infinite range}
\end{itemize}

All dependent on the {\it word size} - maximum size of the string of bits used.

\end{frame}

\begin{frame}{Number systems}
\begin{itemize}
\item{\underline{\bf Decimal:}}
\begin{itemize}
\item{base-10}
\item{digits: 0,1,2,3,4,5,6,7,8,9}
\item{powers of 10 positional system}
\item{Ex: 86409}
\end{itemize}
\end{itemize}
\vspace{3 in}
\end{frame}

\begin{frame}{Number systems}
\begin{itemize}
\item{\underline{\bf Binary:}}
\begin{itemize}
\item{base-2}
\item{digits: 0,1}
\item{powers of 2 positional system}
\item{Ex: 101011}
\end{itemize}
\end{itemize}
\vspace{3 in}
\end{frame}

\begin{frame}{Dec-Bin Conversion}
\begin{itemize}
\item{Convert $185_{10}$ to binary}
\end{itemize}
\vspace{3 in}
\end{frame}

\begin{frame}{Computer representation}
\begin{itemize}
\item{Positive integers}
\begin{itemize}
\item{What we can represent depends on the word size}
\item{Ex: 8-bits, then $43_{10}$ = $00101011_{2}$}
\item{What is the range? $2^8=256$ values from 0 to $2^8 - 1$}
\item{Ex: 16-bits, then $43_{10}$ = $0000000000101011_{2}$}
\item{What is the range? $2^{16}=65,536$ values from 0 to $2^{16} - 1$}
\end{itemize}
\end{itemize}
\end{frame}

\begin{frame}{Computer representation}
\begin{itemize}
\item{Negative integers}
\begin{itemize}
\item{Signed magnitude method - use leftmost bit for the sign}
\item{usually 0 for '+' and 1 for '-'}
\item{Ex: 8-bits, then $-43_{10}$ = $10101011_{2}$}
\item{What is the range? only 7 bits so 128 values with a lead 0 and 128 with a lead 1}
\item{But, two of them - 10000000 and 00000000 - represent the same number, 0}
\item{We usually let 10000000 = -128 giving the range -128 to 127}
\end{itemize}
\end{itemize}
\end{frame}

\begin{frame}{Real numbers - Decimal}
\begin{itemize}
\item{How do we interpet real numbers in decimal?}
\item{Ex: Consider 82.3801}
\vspace{3 in}
\end{itemize}
\end{frame}

\begin{frame}{Real numbers - Binary}
\begin{itemize}
\item{What about in binary?}
\item{Ex: Consider 101.1101}
\vspace{3 in}
\item{Going from decimal to binary with real numbers?}
\vspace{1 in}
\item{Now, how do we represent them in a computer?}
\end{itemize}
\end{frame}

\begin{frame}{Real numbers - Conversion}
\begin{itemize}
\item{Going from decimal to binary with real numbers?}
\vspace{2 in}
\item{Now, how do we represent them in a computer?}
\end{itemize}
\end{frame}


\section{Floating-point numbers}

\begin{frame}{Floating-point number system} 
\begin{itemize}

\item{A floating-point number system is a finite approximation to the
(infinite) real/complex number system, of the form:

\begin{equation} 
\pm m \times b ^ e 
\end{equation}} 

\item{The $m$ is called the {\it \bf mantissa}, $b$ is the {\it \bf base} and
$e$ is the {\it \bf exponent}.}

\end{itemize}
\vspace{1 in}
\end{frame}

\begin{frame}{Floating-point representation} 
\begin{itemize}

\item{How do we store a floating-point number in a word?}
\item{We take the component parts of the floating-point number, (sign, mantissa, exponent) and assign them to different sections of the word}
\item{For example, take $\pm m \times b ^ e$ and put them into the word as follows:} 
\end{itemize}

\begin{center}
\begin{tabular}{ |c|c|c|c| } 
 \hline
 $s_m$ & $s_e$ & exponent & \hspace{0.25 in} mantissa \hspace{0.25 in} \\ 
 \hline
\end{tabular}
\end{center}
\vspace{1 in}
\end{frame}


\begin{frame}{Normalized floating-point number system} 
\begin{itemize}

\item{In normalized floating-point number systems real numbers are represented in the form: 

\begin{equation} 
\pm 0.d_1d_2d_3 \dots d_k \times b ^ e 
\end{equation}}

\item{Where the first digit of the {\it mantissa} should be non-zero, i.e. $ 1
\leq d_1 \leq b-1$.} 
\item{The remaining digits can be zero and are also
constrained by the base, i.e. $ 0 \leq d_i \leq b-1$.}

\item{Because $d_1 \neq 0$, $k$ is the number of significant digits in the mantissa, and is called
the {\it precision} of the floating point system.} 

\end{itemize}
\end{frame}

\begin{frame}{Example - IEEE 754}
\begin{itemize}
\item{binary32 is the single-precision floating point representation}
\item{it uses a 32-bit word with 1 bit for the sign of the mantissa, $s_m$, 8 bits for the signed (biased by 127) exponent, and the remaining 23 bits for the mantissa}
\end{itemize}

\begin{center}
\begin{tabular}{ |c|c|c| } 
 \hline
 $s_m$ & exponent & \hspace{0.65 in} mantissa \hspace{0.65 in} \\ 
 \hline
\end{tabular}
\begin{tabular}{ ccc } 
 $_{1 bit}$ & \hspace{0.15 in} $_{8 bits}$ \hspace{0.28 in} & \hspace{0.65 in} $_{23 bits}$ \hspace{0.65 in} \\ 
\end{tabular}
\end{center}
\vspace{1 in}
\begin{itemize}
\item{this sytem has a precision of $k=24$, don't need to store the lead 1}
\end{itemize}
\end{frame}

\begin{frame}{Errors in floating-point representation}

There are a number of inherent errors in this system, some more obvious than others.

\begin{itemize}
\item{Large negative and positive numbers fall outside the finite range of the system (overflow).} 
\item{Because of normalization, very small  (close to 0) negative and positive numbers fall outside the range (underflow).} 
\item{Only a finite number of values can be represented in the range (round-off error)} 
\item{The distance between two consecutive floating-point numbers increases as the numbers get larger}
\end{itemize}
\end{frame} 

\begin{frame}{Ex: Hypothetical Floating-Point Computer}
Suppose your computer uses a 7-bit word to represent normalized floating-point numbers as follows:

\begin{center}
\begin{tabular}{|c|c|c|c|c|c|c|}
\hline 
$s_m$ & $s_e$ & $e_1$ & $e_2$ & $b_1$ & $b_2$ & $b_3$ \\ 
\hline 
\end{tabular}
\end{center}

where $s_m$ is the sign of the mantissa and  $s_e$ is the sign of the exponent (0 for positive, 1 for negative), 
$b_1,b_2,$ and $b_3$ are the bits of the mantissa, and $e_1,e_2$ are the bits 
of the exponent.
\end{frame}

\begin{frame}{Example continued}
\begin{enumerate}
\item[(a)] What is the smallest positive non-zero number that can be represented in this system? What is it's value in decimal?
\vspace{0.1 in}
\item[(b)] What is the next smallest positive non-zero number that can be represented in this system? 
\vspace{0.1 in}
\item[(c)] What is the distance between the two values above? What does this value represent?
\vspace{0.1 in}
\item[(d)] Use (c) to predict the next value in decimal. Convert it to floating-point to test your answer.
\vspace{0.1 in}
\end{enumerate} 

\end{frame} 




\end{document}

